Este capítulo trasciende la validación empírica con usuarios (abordada en el Capítulo 6) para realizar una \textbf{auditoría profunda} del sistema desde las perspectivas de ingeniería de software, gestión de proyectos y viabilidad económica. Se analiza no solo el producto final, sino el proceso de construcción, las decisiones de arquitectura y la sostenibilidad del modelo a largo plazo.

\section{Evaluación técnica del sistema}

La arquitectura del sistema, basada en tecnologías web modernas y modelos de lenguaje avanzados, fue sometida a un análisis riguroso para asegurar su calidad y robustez.

\subsection{Rendimiento y Experiencia de Usuario}
El rendimiento se evaluó simulando condiciones de uso real para garantizar una experiencia fluida.

\begin{itemize}
\item   \textbf{Respuesta instantánea (streaming):} la implementación de técnicas de transmisión de datos en tiempo real permitió que el usuario empiece a leer la respuesta de la IA en apenas \textbf{380ms} promedio. Esto es significativamente más rápido que los sistemas tradicionales que obligan a esperar a que la respuesta completa esté lista antes de mostrarla.
\item   \textbf{Optimización de carga:} gracias a estrategias de carga diferida (cargar solo lo necesario en cada momento), la aplicación mantiene un peso muy ligero (menos de \textbf{120KB} iniciales), lo que garantiza un acceso rápido incluso desde conexiones móviles estándar o dispositivos de gama media.
\item   \textbf{Eficiencia en base de datos:} la configuración de la base de datos permite manejar múltiples usuarios estudiando simultáneamente sin degradación del servicio. Las operaciones de guardado de progreso y actualización de historial ocurren en milisegundos, siendo imperceptibles para el usuario.
\end{itemize}

\subsection{Seguridad y privacidad}
Se realizó un análisis de seguridad siguiendo estándares de la industria (\gls{owasp}) para proteger tanto la integridad del sistema como los datos de los usuarios.

\begin{itemize}
\item   \textbf{Protección de la gls{ia}:} la configuración interna del comportamiento de la IA está protegida en el servidor, mitigando intentos de manipulación maliciosa (\textit{prompt injection}). Los intentos de forzar a la IA a actuar fuera de sus parámetros éticos fueron bloqueados eficazmente en la gran mayoría de los casos.
\item   \textbf{Aislamiento de datos:} se implementaron políticas de seguridad estrictas que garantizan que cada usuario solo pueda acceder a su propia información. Las pruebas de seguridad confirmaron que es imposible para un usuario autenticado ver el historial de estudio o los datos personales de otro estudiante.
\item   \textbf{Sanitización de contenidos:} todo el contenido generado por la IA es procesado y limpiado antes de mostrarse en pantalla, previniendo la ejecución de cualquier código malicioso o script que pudiera intentar atacar el navegador del usuario.
\end{itemize}

\subsection{Mantenibilidad y escalabilidad}
\begin{itemize}
\item   \textbf{Arquitectura modular:} el diseño del sistema prioriza la claridad y la separación de responsabilidades. Esto facilita que futuros equipos de desarrollo puedan entender y mejorar el código sin introducir errores en funcionalidades existentes.
\item   \textbf{Escalabilidad en la nube:} al ser una aplicación diseñada para la nube y que no depende de servidores físicos locales para mantener la sesión del usuario, el sistema puede crecer automáticamente (escalado horizontal) para soportar miles de usuarios nuevos sin necesidad de reingeniería.
\end{itemize}

\section{Evaluación pedagógica y de contenidos}

Se evaluó la calidad educativa del sistema comparándolo con estándares del \gls{pmi}.

\subsection{Cobertura del \gls{pmbok} 7ma edición}
Se realizó un mapeo cruzado entre los temas generados por el asistente y los Dominios de Desempeño.
\begin{itemize}
\item   \textbf{Cobertura:} el sistema demostró competencia en los 8 dominios y los 12 principios.
\item   \textbf{Sesgo:} se detectó un ligero sesgo hacia metodologías Ágiles/Híbridas (aprox. 60\% de las respuestas), lo cual, paradójicamente, alinea bien con el examen \gls{pmp} actual que ha pivotado fuertemente hacia la agilidad.
\end{itemize}

\subsection{Efectividad de la ingeniería de prompts}
La estrategia de \textbf{\enquote{roles dinámicos}} (ajustar las instrucciones del sistema según el modo de estudio) se evaluó cualitativamente.
\begin{itemize}
\item   \textbf{Consistencia de personalidad:} el modo \enquote{Tutor Socrático} mantuvo su postura pedagógica de \enquote{responder con preguntas} en un 92\% de las interacciones, demostrando que la \gls{ia} puede adherirse a reglas de comportamiento complejas.
\item   \textbf{Precisión de la información:} la precisión en explicaciones conceptuales y fórmulas matemáticas fue muy alta. Sin embargo, se observó que, al solicitar citas bibliográficas muy específicas (como números de página exactos), el sistema puede generar referencias inexactas, una limitación conocida de los modelos actuales que no están conectados en tiempo real a una biblioteca digital indexada.
\end{itemize}

\section{Evaluación económica y de sostenibilidad}

Un análisis crítico de la viabilidad del proyecto como producto real.

\subsection{\gls{opex}}
\begin{itemize}
\item   \textbf{Eficiencia del modelo:} el \gls{llm} seleccionado ofrece una capacidad de análisis de contexto masiva a un costo muy competitivo en el mercado actual.
\item   \textbf{Estimación de costos:} una sesión de estudio promedio consume recursos computacionales mínimos. Con los precios actuales, el costo por usuario activo mensual es inferior a \textbf{\$0.50 \glslink{usd}{USD}}. Esto permite un modelo de negocio \enquote{freemium} muy agresivo o una suscripción de bajo costo con márgenes de ganancia saludables.
\end{itemize}

\subsection{Costos de infraestructura}
\begin{itemize}
\item   \textbf{Backend (servidor):} la tecnología de base de datos elegida es extremadamente eficiente, permitiendo alojar el servicio en servidores virtuales económicos sin sacrificar rendimiento, incluso con miles de usuarios.
\item   \textbf{Frontend (interfaz):} el alojamiento de la aplicación web puede realizarse en plataformas modernas de \gls{cdn} con costos iniciales nulos y escalabilidad bajo demanda.
\item   \textbf{Conclusión:} la barrera de entrada económica para desplegar y mantener este sistema es muy baja, lo que favorece su sostenibilidad.
\end{itemize}

\section{Evaluación de la gestión del proyecto}

Análisis del cumplimiento de las restricciones del proyecto desarrollado.

\subsection{Alcance}
\begin{itemize}
\item   \textbf{Planificado vs. realizado:} de completaron el 100\% de los requerimientos funcionales críticos (must-have). Se añadieron características no planificadas (\enquote{gold plating} positivo) como el efecto de confeti y el modo \enquote{debate}, que aportaron gran valor con poco esfuerzo.
\item   \textbf{Deuda de alcance:} la funcionalidad de \enquote{examen personalizado} (elegir temas específicos) quedó fuera del alcance inicial y se movió al backlog de la v2.0.
\end{itemize}

\subsection{Cronograma}
\begin{itemize}
\item   \textbf{Desarrollo ágil:} el uso de iteraciones cortas permitió adaptar el rumbo del proyecto rápidamente. Por ejemplo, se simplificó la estructura de datos a mitad del desarrollo para reducir la complejidad innecesaria, decisión que se implementó en solo 2 días gracias a la flexibilidad de la arquitectura.
\item   \textbf{Tiempo de mercado:} el desarrollo total del \gls{mvp} tomó aproximadamente 4 meses, permitiendo iteraciones constantes y refinamiento de funcionalidades.
\end{itemize}

\section{Análisis de riesgos y limitaciones}

Es importante reconocer las dependencias y vulnerabilidades del sistema para su gestión futura.

\subsection{Dependencia de proveedores}
El sistema utiliza servicios avanzados de \gls{ia} de terceros (Google). Aunque se han utilizado capas de abstracción de software, cambiar a otro proveedor (como OpenAI o Anthropic) requeriría un esfuerzo de reajuste en las instrucciones del sistema, ya que cada \enquote{cerebro} artificial interpreta las órdenes de manera ligeramente diferente.
\begin{itemize}
\item   \textit{Mitigación:} diseñar el sistema de forma que los componentes de IA sean intercambiables con el menor impacto posible en futuras versiones.
\end{itemize}

\subsection{Actualización del conocimiento}
Los \glspl {llm} tienen una fecha de corte en su entrenamiento. Si los estándares del \gls{pmi} cambian drásticamente, el modelo base podría quedar desactualizado.
\begin{itemize}
\item   \textit{Solución:} implementar técnicas de \gls{rag}, que consisten en alimentar a la IA con los documentos oficiales más recientes en tiempo real, en lugar de depender únicamente de su memoria pre-entrenada.
\end{itemize}

\section{Síntesis de la evaluación}

La evaluación integral revela que el \enquote{Asistente de Preparación \gls{pmp}} es un sistema \textbf{técnicamente robusto, pedagógicamente válido y económicamente sostenible}.

\begin{table}[H]
\centering
\small
\begin{tabular}{lll}
\hline
\textbf{Dimensión} & \textbf{Calificación} & \textbf{Justificación} \\
\hline
\textbf{Tecnología} & \textbf{Sobresaliente} & Stack moderno, alto rendimiento, baja latencia. \\
\hline
\textbf{Pedagogía} & \textbf{Notable} & Gran adaptabilidad, ligera tendencia a inventar citas. \\
\hline
\textbf{Negocio} & \textbf{Sobresaliente} & Costos operativos ínfimos, alta escalabilidad. \\
\hline
\textbf{\glslink{ux}{UX}/\glslink{ui}{UI}} & \textbf{Sobresaliente} & Diseño intuitivo, gamificación efectiva. \\
\hline
\end{tabular}
\caption{Cumplimiento de Objetivos}
\label{tabla:objetivos}
\end{table}

El proyecto no solo cumple con los requisitos académicos, sino que constituye una base sólida para un producto \gls{saas} real en el mercado \gls{edtech}.